\documentclass[11pt]{article}
\usepackage[margin=1in]{geometry}
\usepackage{amsmath,amssymb,amsthm,mathtools}
\usepackage{booktabs,array,hyperref,microtype}
\usepackage[T1]{fontenc}
\usepackage{lmodern}
\newtheorem{theorem}{Theorem}[section]
\newtheorem{proposition}[theorem]{Proposition}
\newtheorem{lemma}[theorem]{Lemma}
\newtheorem{corollary}[theorem]{Corollary}
\newtheorem{definition}[theorem]{Definition}
\newtheorem{remark}[theorem]{Remark}
\newcommand{\C}{\mathbb C}
\newcommand{\Z}{\mathbb Z}
\newcommand{\R}{\mathbb R}
\newcommand{\Argp}{\operatorname{Arg}}
\newcommand{\Logp}{\operatorname{Log}}
\newcommand{\wtC}{\widetilde{\C^{\times}}}
\title{Branch-Indexed Multiplication and Logarithm Algebra\\on the Universal Cover of $\C^{\times}$}
\author{Donald G. Palmer}
\date{July 2026}
\begin{document}
\maketitle

\begin{abstract}
We formalize the branch-index layer used in the Complex Numeric Representational System (CNRS) as a concrete model of the universal covering group of the nonzero complex numbers.  A lifted value is represented by a pair $(z,k)$ with $z\in\C^{\times}$ and $k\in\Z$, where the lifted argument is $\Argp z+2\pi k$.  We derive the exact branch-wrap cocycle governing multiplication, prove associativity, commutativity, inversion, and integer-power laws, and show that the lifted logarithm is a single-valued group isomorphism from the covering group to the additive complex plane.  We also distinguish this multiplicative algebra from addition: ordinary complex addition does not induce a canonical globally defined group law on the cover of $\C^{\times}$.  The resulting theorem package gives a precise mathematical foundation for CNRS branch-index multiplication and logarithmic bookkeeping while remaining independent of CNRS-A digit representation.
\end{abstract}

\section{Purpose and status}
CNRS Layer 2 records the sheet of the complex logarithm by adjoining an integer branch index to a nonzero complex value.  The purpose of this note is to replace informal branch bookkeeping by a complete algebraic construction.

The results are standard consequences of covering-space theory, but the explicit pair model, cocycle, and computational formulas are the structures required by CNRS.  The contribution here is therefore an exact alignment theorem between the CNRS branch-index representation and the universal covering group of $\C^{\times}$.

Throughout, the principal argument is
\[
\Argp z\in(-\pi,\pi],
\]
and the principal logarithm is
\[
\Logp z=\ln|z|+i\Argp z.
\]
The point $z=0$ is excluded.

\section{Two equivalent models of the cover}

\begin{definition}[Angle model]
Define
\[
\mathcal U=\R\times\R
\]
with group law
\[
(u,\theta)+(v,\phi)=(u+v,\theta+\phi)
\]
and covering map
\[
p(u,\theta)=e^{u+i\theta}\in\C^{\times}.
\]
\end{definition}

This is the standard simply connected universal cover of $\C^{\times}$.

\begin{definition}[Branch-index model]
Define
\[
\wtC=\C^{\times}\times\Z.
\]
For $(z,k)\in\wtC$, its lifted argument and lifted logarithm are
\[
\Theta(z,k)=\Argp z+2\pi k,
\qquad
\widetilde{\operatorname{Log}}(z,k)=\ln|z|+i\Theta(z,k).
\]
\end{definition}

\begin{proposition}
The map
\[
F:\wtC\longrightarrow\mathcal U,
\qquad
F(z,k)=\bigl(\ln|z|,\Argp z+2\pi k\bigr)
\]
is a bijection.
\end{proposition}

\begin{proof}
Injectivity follows because equality of the first coordinates gives equal moduli, while equality of lifted angles gives equal arguments modulo $2\pi$ and hence equal complex values; the principal-argument convention then forces equal branch indices.  For surjectivity, given $(u,\theta)$ let $z=e^{u+i\theta}$ and let
\[
k=\frac{\theta-\Argp z}{2\pi}\in\Z.
\]
Then $F(z,k)=(u,\theta)$.
\end{proof}

\section{The branch-wrap cocycle}

\begin{definition}[Wrap cocycle]
For $z,w\in\C^{\times}$ define
\[
\omega(z,w)=\frac{\Argp z+\Argp w-\Argp(zw)}{2\pi}.
\]
\end{definition}

\begin{lemma}
For all $z,w\in\C^{\times}$, $\omega(z,w)\in\Z$.  With the convention $\Argp\in(-\pi,\pi]$, one has
\[
\omega(z,w)\in\{-1,0,1\}.
\]
\end{lemma}

\begin{proof}
The two quantities $\Argp z+\Argp w$ and $\Argp(zw)$ represent the same angle modulo $2\pi$, so their difference is an integer multiple of $2\pi$.  Since the first lies in $(-2\pi,2\pi]$ and the second in $(-\pi,\pi]$, only the three displayed integers can occur.
\end{proof}

An explicit rule is
\[
\omega(z,w)=
\begin{cases}
 1,&\Argp z+\Argp w>\pi,\\
-1,&\Argp z+\Argp w\le -\pi,\\
 0,&\text{otherwise},
\end{cases}
\]
where the lower boundary reflects the half-open principal interval.

\begin{theorem}[Cocycle identity]
For all $z,w,u\in\C^{\times}$,
\[
\omega(z,w)+\omega(zw,u)
=
\omega(w,u)+\omega(z,wu).
\]
\end{theorem}

\begin{proof}
Both sides equal
\[
\frac{\Argp z+\Argp w+\Argp u-\Argp(zwu)}{2\pi}
\]
after cancellation of the intermediate principal arguments.
\end{proof}

\section{Lifted multiplication}

\begin{definition}[Branch-index multiplication]
For $(z,k),(w,\ell)\in\wtC$, define
\[
(z,k)\odot(w,\ell)
=
\bigl(zw,\,k+\ell+\omega(z,w)\bigr).
\]
\end{definition}

\begin{theorem}[Covering-group algebra]
The set $(\wtC,\odot)$ is an abelian group.  Its identity is $(1,0)$, and the projection
\[
\pi:\wtC\to\C^{\times},\qquad \pi(z,k)=z
\]
is a surjective group homomorphism with kernel
\[
\ker\pi=\{(1,k):k\in\Z\}\cong\Z.
\]
\end{theorem}

\begin{proof}
Closure is immediate.  Associativity follows from the cocycle identity.  Commutativity follows from $zw=wz$ and symmetry of $\omega$.  Since $\Argp 1=0$, $(1,0)$ is an identity.  The projection is visibly a homomorphism and is surjective.  Its kernel is exactly the displayed set.
\end{proof}

\begin{proposition}[Inverse]
Let $x=(z,k)$.  Then
\[
x^{-1}=
\left(z^{-1},-k-\omega(z,z^{-1})\right).
\]
Equivalently, if $\Argp z\ne\pi$ then
\[
x^{-1}=(z^{-1},-k),
\]
while for negative real $z$ (so $\Argp z=\pi$),
\[
x^{-1}=(z^{-1},-k-1).
\]
\end{proposition}

\begin{proof}
The formula is forced by requiring the second component of $x\odot x^{-1}$ to vanish.  The boundary case occurs because $\Argp z=\Argp z^{-1}=\pi$, so $\omega(z,z^{-1})=1$.
\end{proof}

\begin{proposition}[Integer powers]
For $n\in\Z$, the power $x^{\odot n}$ is uniquely defined by the group law and satisfies
\[
\Theta(x^{\odot n})=n\Theta(x),
\qquad
\pi(x^{\odot n})=z^n.
\]
Consequently
\[
x^{\odot n}
=
\left(z^n,\frac{n(\Argp z+2\pi k)-\Argp(z^n)}{2\pi}\right).
\]
\end{proposition}

\section{The single-valued lifted logarithm}

\begin{theorem}[Logarithm isomorphism]
The lifted logarithm
\[
\widetilde{\operatorname{Log}}: (\wtC,\odot)\longrightarrow(\C,+),
\qquad
\widetilde{\operatorname{Log}}(z,k)=\ln|z|+i(\Argp z+2\pi k),
\]
is a group isomorphism.
\end{theorem}

\begin{proof}
For $x=(z,k)$ and $y=(w,\ell)$,
\begin{align*}
\widetilde{\operatorname{Log}}(x\odot y)
&=\ln|zw|+i\left(\Argp(zw)+2\pi(k+\ell+\omega(z,w))\right)\\
&=\ln|z|+\ln|w|
+i\left(\Argp z+\Argp w+2\pi k+2\pi\ell\right)\\
&=\widetilde{\operatorname{Log}}(x)+\widetilde{\operatorname{Log}}(y).
\end{align*}
Bijectivity follows from the earlier bijection $F$ and the identification $(u,\theta)\mapsto u+i\theta$.
\end{proof}

\begin{corollary}[Exact product law]
For all $x,y\in\wtC$,
\[
\widetilde{\operatorname{Log}}(x\odot y)
=
\widetilde{\operatorname{Log}}(x)+\widetilde{\operatorname{Log}}(y)
\]
without an unspecified $2\pi i n$ term.
\end{corollary}

\begin{corollary}[Lifted exponential]
Define
\[
\widetilde{\operatorname{Exp}}(a+ib)
=
\left(e^{a+ib},\frac{b-\Argp(e^{a+ib})}{2\pi}\right).
\]
Then
\[
\widetilde{\operatorname{Log}}\circ\widetilde{\operatorname{Exp}}=\operatorname{id}_{\C},
\qquad
\widetilde{\operatorname{Exp}}\circ\widetilde{\operatorname{Log}}=\operatorname{id}_{\wtC}.
\]
\end{corollary}

Thus the familiar many-valued logarithm is replaced on the cover by a globally single-valued inverse of the exponential map.

\section{Fractional powers}
For $x\in\wtC$ and $a\in\C$, define the lifted power
\[
x^{[a]}=\widetilde{\operatorname{Exp}}\bigl(a\,\widetilde{\operatorname{Log}}x\bigr).
\]
This is single-valued on the cover.  It obeys
\[
x^{[a+b]}=x^{[a]}\odot x^{[b]},
\]
and for fixed $a$,
\[
(x\odot y)^{[a]}
=
\widetilde{\operatorname{Exp}}\left(a\widetilde{\operatorname{Log}}x+a\widetilde{\operatorname{Log}}y\right).
\]
When $a\in\Z$, this agrees with group powers.  Projecting to $\C^{\times}$ recovers one branch-selected value of the usual multivalued power.

\section{Why addition is different}
Multiplication preserves $\C^{\times}$ and is the group operation whose universal cover is being constructed.  Addition has three obstructions.

First, $z+w$ may be zero, which lies outside $\C^{\times}$.  Second, a branch index records winding around zero under multiplication; it carries no canonical information about a path from $z$ and $w$ to $z+w$.  Third, any rule
\[
(z,k)\boxplus(w,\ell)=(z+w,m)
\]
requires a choice of $m$ that is not determined by the endpoint $z+w$ and the two individual sheet indices alone in a way compatible with arbitrary continuous deformations through the punctured plane.

\begin{proposition}[No global lifted additive group]
There is no everywhere-defined binary operation $\boxplus$ on $\wtC$ for which:
\begin{enumerate}
\item $\pi(x\boxplus y)=\pi(x)+\pi(y)$ for all $x,y$;
\item $(\wtC,\boxplus)$ is a group;
\item $\pi$ maps into $\C^{\times}$.
\end{enumerate}
\end{proposition}

\begin{proof}
Choose $y$ with $\pi(y)=-\pi(x)$.  Then condition 1 requires $\pi(x\boxplus y)=0$, contradicting the codomain $\C^{\times}$.
\end{proof}

A partial lifted addition can be defined on pairs with $z+w\ne0$ after selecting a path or a branch convention, but it is not intrinsic to the universal covering group and should not be conflated with the multiplication law proved above.

\section{CNRS interpretation}
The branch-index layer can now be stated precisely:
\begin{itemize}
\item CNRS-A may represent the projected value $z$ numerically.
\item The integer $k$ records the logarithmic sheet.
\item Multiplication combines the projected values and updates $k$ by the wrap cocycle.
\item The lifted logarithm is single-valued and additive.
\item Inversion, powers, and roots are exact operations on the cover.
\item Addition remains ordinary complex addition at the projected level and requires separate branch/path handling if a lifted result is desired.
\end{itemize}

The construction does not depend on the base $-2+i$ or on any unresolved CNRS completeness theorem.  It is a general topological-algebraic layer that can be paired with any exact or approximate representation of nonzero complex values.

\section{Implementation specification}
A computational implementation should store
\[
(z,k),\qquad z\ne0,
\]
with a single locked principal-argument convention.  It should expose:
\begin{verbatim}
wrap_index(z, w)
lifted_multiply((z, k), (w, l))
lifted_inverse((z, k))
lifted_power_integer((z, k), n)
lifted_log((z, k))
lifted_exp(value)
lifted_power((z, k), exponent)
\end{verbatim}
The implementation must test the branch-cut boundary, especially negative real inputs, because $\Argp(-r)=\pi$ rather than $-\pi$ under the convention used here.

\section{Claim status and open questions}
The covering-group, cocycle, inverse, power, and logarithm results are established mathematical results in the explicit CNRS pair model.  They complete the multiplicative/logarithmic algebra of the branch-index layer.

The following remain separate questions:
\begin{enumerate}
\item integration of the pair model with canonical CNRS-A representations;
\item exact branch-aware addition under specified path data;
\item analytic continuation of CNRS-H functions on the cover;
\item extension to matrices, operators, or fields with zeros and singularities;
\item physical interpretation of branch indices in Scale Space applications.
\end{enumerate}

\section*{Conclusion}
The CNRS branch index is not merely an auxiliary label.  With the wrap cocycle it gives an exact coordinate model of the universal covering group of $\C^{\times}$.  On this cover, multiplication is globally well defined, the logarithm is single-valued and additive, and inversion and powers have exact branch-aware formulas.  This closes the core Layer-2 multiplication/logarithm problem while showing precisely why addition belongs to a different structure.

\section*{References}
\begin{enumerate}
\item L. V. Ahlfors, \emph{Complex Analysis}, 3rd ed., McGraw--Hill, 1979.
\item O. Forster, \emph{Lectures on Riemann Surfaces}, Springer, 1981.
\item A. Hatcher, \emph{Algebraic Topology}, Cambridge University Press, 2002.
\item A. E. Bashirov and S. Norozpour, ``Riemann surface of complex logarithm and multiplicative calculus,'' arXiv:1610.00133, 2016.
\end{enumerate}

\section*{AI assistance declaration}
Artificial-intelligence tools were used to organize the derivation, check formulas, prepare computational verification, and draft the manuscript.  The author remains responsible for the mathematical claims, interpretation, and final text.

\end{document}
